\title{Direct Preference Optimization: Your Language Model is Secretly a Reward Model}
We may now apply this formulation to the ground-truth reward $r^*$ and its associated optimal policy $\pi^*$. Importantly, the Bradley-Terry model requires only the difference between the rewards assigned to two candidate completions, i.e., ${p^*(y_1 \succ y_2 \mid x) = \sigma(r^*(x, y_1) - r^*(x, y_2))}$. By substituting the expression from Eq.~\ref{eq:main_eq} for $r^*(x,y)$ into the preference probability equation (Eq.~\ref{eq:bradley-terry}), the partition function terms cancel, which allows the preference probability to be represented solely in terms of the optimal policy $\pi^*$ and the reference policy $\pi_\text{ref}$. Thus, under the Bradley-Terry framework, the optimal RLHF policy $\pi^*$ adheres to the following preference structure:

We provide a summary of the RLHF methodology as described by \citeauthor{ziegler2020finetuning} and further developed in \citep{stiennon2022learning, bai2022training, ouyang2022training}. The standard workflow consists of three main components: (1) supervised fine-tuning (SFT), (2) collection of preference data and training of the reward model, and (3) reinforcement learning-based policy optimization.

RLHF algorithms generally maximize a reward function while imposing a KL-divergence constraint to prevent the policy from straying excessively from a reference policy. Thus, in evaluating these methods, it is crucial to consider not only the attained reward but also the associated KL divergence; a marginal increase in reward at the cost of a much larger KL may not be preferable. Figure~\ref{fig:frontier-tldr-main} presents the trade-off between reward and KL for different approaches within the sentiment analysis context. For each algorithm, we conduct multiple experiments, varying the hyperparameters that control the degree of policy conservativeness (for PPO, target KL $\in\{3,6,9,12\}$; for DPO, $\beta \in \{0.05,0.1,1,5\}$; for unlikelihood, $\alpha\in\{0.05,0.1,0.5,1\}$; for preferred-FT, we use distinct random seeds), for a total of 22 experimental runs. Policies are evaluated every 100 training steps until they converge, assessing their mean reward using the true reward function as well as the average sequence-level KL divergence\footnote{Summing the KL-divergences computed at each time step.} relative to the reference policy, $\text{KL}\left(\pi\mid \mid \pi_\text{ref}\right)$. The results show that DPO delivers a substantially more favorable reward-KL tradeoff compared to alternatives, achieving top rewards while maintaining low KL values. This finding is significant for several reasons. Most notably, although DPO and PPO share the same objective, DPO demonstrates markedly superior efficiency—its reward/KL frontier universally outperforms that of PPO. Furthermore, DPO surpasses PPO’s frontier even when PPO has access to ground truth rewards (PPO-GT).

\begin{equation}\label{eq:optimum_model}
    \mathcal{L}_\text{DPO}(\pi_{\theta}; \pi_\text{ref}) = -\mathbb{E}_{(x, y_w, y_l)\sim \mathcal{D}}\left[\log \sigma \left(\beta \log \frac{\pi_{\theta}(y_w\mid x)}{\pi_\text{ref}(y_w\mid x)} - \beta \log \frac{\pi_{\theta}(y_l\mid x)}{\pi_\text{ref}(y_l\mid x)}\right)\right].
\end{equation}

\textbf{Tasks.} Our study examines three distinct open-ended text generation challenges. In every case, the learning algorithm infers a policy using a preference dataset $\mathcal{D}=\bigl\{x^{(i)}, y_w^{(i)}, y_l^{(i)}\bigr\}_{i=1}^N$. For the \textbf{controlled sentiment generation} task, $x$ is provided as an initial segment from a movie review sourced from the IMDb dataset \cite{maas-EtAl:2011:ACL-HLT2011}, and the objective is to have the policy generate a continuation $y$ exhibiting positive sentiment. To facilitate a systematic assessment, for this task we automatically construct preference pairs among generations utilizing a pre-trained sentiment classifier, such that $p(\text{positive}\mid x,y_w)>p(\text{positive}\mid x,y_l)$. For SFT, GPT-2-large is fine-tuned on the IMDB training set reviews until performance plateaus (additional specifics are included in App~\ref{app:sentiment_details}). In the case of \textbf{summarization}, $x$ represents a Reddit forum post, and the policy aims to compose a summary $y$ covering the post’s main ideas. Consistent with existing literature, we use the Reddit TL;DR summarization dataset \citep{volske-etal-2017-tl}, supplemented by human-labeled preferences collected by \citeauthor{stiennon2022learning}. Our SFT model is adapted on human-generated Reddit forum post summaries\footnote{\url{https://huggingface.co/CarperAI/openai_summarize_tldr_sft}}, employing the TRLX framework \citep{leandro_von_werra_2023_7790115} for RLHF. Note that the human preference data originates from the work of \citeauthor{stiennon2022learning} and consists of outputs from another similarly-trained SFT model. Lastly, for the \textbf{single-turn dialogue} task, $x$ refers to a user-issued query, which could range from scientific questions to personal advice requests. Here, the policy is tasked with generating a helpful, engaging reply $y$ for the user; we leverage the Anthropic Helpful and Harmless dialogue dataset \citep{bai2022training}, which comprises 170,000 dialogues between human participants and an automated assistant. Each conversation concludes with two responses, both produced by a sizable (but unspecified) language model, along with an annotation indicating which one was preferred by the human. Since there is no available pre-trained SFT model in this case, we construct the SFT model by fine-tuning a generic language model exclusively on the preferred responses.

\subsection{Instability of Actor-Critic Algorithms} Our framework also offers a lens through which to analyze the instabilities observed in standard actor-critic algorithms deployed for RLHF, including PPO. Focusing on the RL fine-tuning phase delineated in Section \ref{section:prelims}, and drawing parallels to the control-as-inference paradigm \cite{levine2018reinforcement} for the constrained RL formulation in \ref{eq:RL}, we assume a parameterized policy $\pi_{\theta}(y\mid x)$ and seek to minimize the divergence $\mathbb{D}_{\text{KL}}[\pi_{\theta}(y|x) \mid \mid \pi^*(y\mid x)]$, where $\pi^*$ denotes the optimal policy from Eq. \ref{eq:optimum_model}, induced by the reward function $r_{\phi}(y, x)$. This leads, through some algebraic manipulation, to the following objective:

\textbf{RL Fine-Tuning Phase}: In the subsequent reinforcement learning phase, the policy leverages the trained reward function to guide language model adaptation. Consistent with prior works~\citep{jaques2017sequence, jaques2020human}, the objective is formulated as

\begin{equation}\label{eq:bradley-terry}
    p^*(y_1\succ y_2 \mid x)=\frac{\exp\left(r^*(x, y_1)\right)}{\exp\left(r^*(x, y_1)\right) + \exp\left(r^*(x, y_2)\right)}.
\end{equation}

\textbf{Limitations \& Future Directions.} Our findings prompt several avenues for further investigation. One critical question is how well DPO-trained policies extrapolate to domains or distributions not seen during training, particularly in comparison to approaches that learn from explicit reward functions. Preliminary evidence indicates that DPO may achieve generalization comparable to PPO-based strategies, but more systematic evaluation is necessary. For instance, it remains to be seen if leveraging self-labeled data from the DPO policy can be as effective for harnessing unannotated prompts. Another open question pertains to the manifestation of reward over-optimization within the direct preference optimization paradigm—specifically, whether the slight performance decline observed in Figure~\ref{fig:dialogue-main}-right is indicative of this phenomenon. Furthermore, while our experiments span models up to 6B parameters, scaling DPO to models of vastly larger size is a promising topic for future exploration. With respect to evaluation practices, we observe that GPT-4-based win rates are sensitive to prompt formulation; optimizing automated evaluation protocols thus remains an important area for research. Beyond language modeling, DPO holds potential for a wide array of generative modeling applications across different modalities, expanding the scope of preference-based optimization.

\subsection{Your Language Model Is Secretly a Reward Model} DPO eliminates the need for both constructing a separate reward model and performing reinforcement learning for policy optimization, as it employs a unified maximum likelihood objective. The optimization problem given by Eq. \ref{eq:main_eq} can be reformulated as a Bradley-Terry model where rewards are expressed as $r^*(x, y) = \beta \log\frac{\pi^*_\theta(y \mid x)}{\pi_\text{ref}(y \mid x)}$. We train the parametric model $\pi_{\theta}$, which is equivalent to optimizing a reward model according to Eq. \ref{eq:reward_model} after an appropriate variable transformation. In what follows, we develop the theoretical underpinnings of this reparameterization, demonstrate that it imposes no restrictions on the reward function class that can be represented, and prove that it can recover the optimal policy exactly. We start by introducing an equivalence relation over reward functions.

\textbf{Reward Modelling Phase}: The next stage involves prompting the SFT model with various inputs $x$ to generate answer pairs $(y_1, y_2) \sim \pi^\text{SFT}(y \mid x)$. Human annotators are then asked to indicate their preferred response, recorded as $y_w \succ y_l \mid x$, with $y_w$ representing the favored answer and $y_l$ the less preferred option among $(y_1, y_2)$. These preferences are assumed to originate from some hidden reward function $r^*(y, x)$, which is not directly observable. There are multiple frameworks for modeling such preferences; one widely adopted approach is the Bradley-Terry (BT) model \cite{bradley1952rankanalysis}, though more general ranking models like the Plackett-Luce model \citep{plackett1975analysis, luce2012individual} can also be integrated if multiple ranked outputs are available. The BT model expresses the probability distribution of human preferences $p^*$ as:

\begin{multline*}\label{eq:gradient}
    \nabla_\theta \mathcal{L}_\text{DPO}(\pi_\theta;\pi_\text{ref}) = \\ -\beta\mathbb{E}_{(x, y_w, y_l) \sim \mathcal{D}} \left[\underbrace{\sigma(\hat{r}_\theta(x, y_l) - \hat{r}_\theta (x, y_w))}_{\text{greater impact when the reward is misestimated}}\left[\underbrace{\nabla_\theta\log \pi(y_w \mid x)}_{\text{encourage $y_w$}} - \underbrace{\nabla_\theta\log\pi(y_l \mid x)}_{\text{discourage $y_l$}}\right]\right],
\end{multline*}

where $\beta$ regulates the extent to which the new policy $\pi_\theta$ may diverge from the reference policy $\pi_\text{ref}$, which is generally set as the SFT model $\pi^\text{SFT}$. In practice, $\pi_\theta$ is initialized using $\pi^\text{SFT}$. The inclusion of the KL-divergence penalty plays a crucial role in ensuring that the updated policy remains close to the region of input space where the reward function is reliable, as well as in preserving response diversity and preventing collapse to narrow sets of high-reward outputs. Due to the non-differentiable nature of text generation, this objective is conventionally addressed through reinforcement learning algorithms. The prevalent approach \citep{ziegler2020finetuning, stiennon2022learning, bai2022training, ouyang2022training} constructs the augmented reward $r(x, y) = r_{\phi}(x, y) -\beta (\log \pi_{\theta}(y\mid x) - \log \pi_\text{ref}(y\mid x))$ and applies policy gradient methods such as PPO \cite{schulman2017proximal} for maximization.

In this work, we demonstrate a method for optimizing language models to align with human preferences without the need for explicit reward models or reinforcement learning frameworks. We introduce \textit{{Direct Preference Optimization} (DPO)}, an approach that implicitly targets the same objective found in conventional RLHF techniques—namely, maximizing reward under a KL-divergence penalty—yet offers a more accessible and easily trainable alternative. Conceptually, DPO updates the model by amplifying the log-likelihood ratio of preferred over non-preferred outputs, but it further integrates an adaptive, instance-specific weighting mechanism. This innovation helps avert the model collapse issues observed with straightforward probability ratio-based objectives. DPO, akin to prior algorithms, is grounded in a formal preference model, such as the Bradley-Terry model (\cite{bradley1952rankanalysis}), which assesses the correspondence between a candidate reward function and actual preference observations. Distinct from earlier techniques—which use the preference model to craft a loss for training a reward estimator and then a separate policy optimized for that reward—DPO employs a variable transformation to directly express the preference loss in terms of the policy itself. Given a set of human-annotated preference pairs among model responses, DPO enables policy optimization using a straightforward binary cross-entropy loss, yielding the optimal policy under a reward function that is implicitly tailored to the preference data.

\begin{theorem}\label{thm:main}
    Given mild regularity conditions, every equivalence class of rewards compatible with the Plackett-Luce (and specifically Bradley-Terry) models admits a representation of the form ${r(x, y) = \beta \log \frac{\pi(y\mid x)}{\pi_\text{ref}(y\mid x)}}$ for some policy model $\pi(y\mid x)$ and a chosen reference distribution $\pi_\text{ref}(y \mid x)$.
\end{theorem}

\textbf{What is the effect of the DPO update?} To gain mechanistic insight into DPO, it is instructive to examine the gradient of its loss $\mathcal{L}_\text{DPO}$. The gradient with respect to the model parameters $\theta$ is expressed as:

Our principal contribution is Direct Preference Optimization (DPO): a streamlined, reinforcement learning–free methodology for preference-driven training of language models. Empirical results indicate that DPO matches or surpasses the performance of prevailing methods, including PPO-based RLHF, across tasks such as sentiment control, summarization, and conversational modeling, utilizing models with scales up to 6B parameters.

\subsection{Can DPO scale to real preference datasets?}
\label{sec:dpo-real-datasets}
We proceed to assess the fine-tuning capabilities of DPO on tasks involving summarization and single-turn dialogue. In summarization, it has been noted that standard automated metrics like ROUGE often do not align well with human judgments~\citep{stiennon2022learning}. Previous studies have demonstrated that language models fine-tuned with PPO, based on human preference data, can produce more compelling summaries. Our approach involves generating model outputs on the test partition of the TL;DR summarization dataset, and then calculating the average frequency with which these outputs are preferred over the reference summaries in the test set. For all methods, completions are sampled at a range of temperatures from 0.0 to 1.0, with the resulting win rates illustrated in Figure~\ref{fig:frontier-tldr-main} (right). DPO, PPO, and Preferred-FT are each fine-tuned from an identical GPT-J SFT model\footnote{\url{https://huggingface.co/CarperAI/openai_summarize_tldr_sft}}. Our analysis shows that at temperature 0.0, DPO attains a win rate of about 61\%, surpassing PPO, which achieves roughly 57\% at its optimal temperature of 0.0. Furthermore, DPO records a superior maximum win rate relative to the best of $N$ baseline. It should be highlighted that no extensive hyperparameter tuning for DPO’s $\beta$ was performed, indicating that its potential performance may not be fully realized in these results. DPO is observed to maintain stable performance across different sampling temperatures, whereas PPO’s effectiveness diminishes, eventually aligning with the baseline GPT-J model at higher temperatures. Notably, Preferred-FT fails to offer significant improvements over the SFT baseline. Finally, as discussed in Section~\ref{sec:human-judgments}, a direct human comparison of DPO and PPO revealed that samples from DPO at temperature 0.25 were favored 58\% of the time over PPO’s outputs at temperature 0.

\subsection{Validating GPT-4 judgments with human judgments}
\label{sec:human-judgments}

Beyond linguistic applications, the task of learning policies from preferences has received attention in both bandit and reinforcement learning frameworks, with a variety of solutions proposed. In the contextual bandit domain, learning via preferences or action rankings rather than direct rewards defines the contextual dueling bandit (CDB) scenario \citep{yue2012karmed, dudik2015contextual}. For CDBs, where explicit rewards are unavailable, theoretical developments introduce the concept of a \textit{von Neumann winner}: a policy whose expected probability of outperforming any other policy is no less than 50\% \citep{dudik2015contextual}. Unlike the online provision of preference feedback in CDBs, human preference learning typically operates on a fixed, offline dataset of preference-labeled action pairs \citep{yan2022human}. In a related vein, \textit{preference-based RL} (PbRL) utilizes binary preferences derived from an unknown 'scoring' function as opposed to reward signals \citep{BusaFekete2014, ruiz2023dueling}. PbRL encompasses a range of algorithms—some of which exploit off-policy preference information—but most approaches first construct an explicit estimate of the underlying scoring function (effectively a reward model) before proceeding to optimize it \citep{jain2013learning, BusaFekete2014, christiano2017deep, sadigh2017active, kupcsik2018learning}. In contrast, we propose a unified policy learning method that directly seeks to align the policy with observed preferences in a single optimization stage.

\textbf{SFT}: The process is initiated by further training a pre-trained language model using supervised learning on curated datasets tailored for specific tasks (such as dialogue or summarization), yielding a model denoted $\pi^\text{SFT}$.

In this framework, we infer an implicit reward via an alternative parameterization, where the optimal policy directly corresponds to $\pi_\theta$. Additionally, because the approach aligns with learning a reparametrized Bradley-Terry model, it inherits desirable theoretical attributes, including consistency under appropriate assumptions on the preference data distribution \cite{bong2022generalized}. We elaborate further on the theoretical underpinnings of DPO and draw comparisons to related methods in Section~\ref{sec:theory}.

\begin{equation}\label{eq:RL}
\max_{\pi_{\theta}}  \mathbb{E}_{x\sim \mathcal{D}, y\sim \pi_{\theta}(y \mid x)}\bigl[r_{\phi}(x, y)\bigr] - \beta\mathbb{D}_{\textrm{KL}}\bigl[\pi_{\theta}(y\mid x)\mid \mid \pi_\text{ref}(y\mid x)\bigr],
\end{equation}

The objective under consideration is identical to that optimized in previous studies 
\citep{ziegler2020finetuning, stiennon2022learning, bai2022training, ouyang2022training}, utilizing the DPO-equivalent reward within the reward family $r_{\phi}$. Within this framework, the normalization component in $f(r_{\phi}, \pi_\text{ref}, \beta)$ can be understood as representing the soft value function associated with the reference policy $\pi_\text{ref}$. Although this normalization does not impact the optimality of the solution itself, omitting it can result in the policy gradient exhibiting substantial variance, potentially destabilizing the learning process. One approach to handling this normalization is through the use of a learned value function; however, optimizing such a value function is itself a challenging task. Alternatively, earlier works have addressed normalization by employing a human completion baseline, effectively providing a one-sample Monte Carlo estimate of the normalization constant. By contrast, the DPO reparameterization produces a reward function that inherently eliminates the need for any baseline.

\textbf{DPO overview.} The standard DPO framework consists of the following steps: 1) For each prompt $x$, generate completion samples $y_1, y_2 \sim \pi_\text{ref}(\cdot \mid x)$, annotate these with human preferences, and thus construct an offline preference dataset $\mathcal{D} = \{x^{(i)}, y_w^{(i)}, y_l^{(i)}\}_{i=1}^N$; 2) Train the language model $\pi_\theta$ by minimizing the DPO objective $\mathcal{L}_\text{DPO}$ using the reference policy $\pi_\text{ref}$, dataset $\mathcal{D}$, and chosen $\beta$. 
In applied settings, it is preferable to utilize existing public preference datasets rather than creating new samples and collecting fresh human feedback. Because such datasets are usually generated with $\pi^\text{SFT}$, we set $\pi_\text{ref} = \pi^\text{SFT}$ when this model is available. If, however, $\pi^\text{SFT}$ is not provided, $\pi_\text{ref}$ is initialized by maximizing the likelihood of preferred outputs ${(x, y_w)}$; specifically, $\pi_\text{ref} = \argmax_{\pi}\mathbb{E}_{x, y_w \sim \mathcal{D}}\left[\log \pi(y_w \mid x)\right]$. This initialization strategy helps address the distributional gap between the true (but inaccessible) reference distribution and the practical $\pi_\text{ref}$ adopted by DPO. For additional implementation specifics and hyperparameter settings, see Appendix~\ref{app:implementation}.

Theorem~\ref{thm:main} can alternatively be interpreted as pinpointing the specific reward function chosen by the DPO reparameterization within each equivalence class—namely, the function that fulfills:

\subsection{How well can DPO optimize the RLHF objective?}

The arguments supporting these statements are elementary, so we postpone their formal demonstration to Appendix \ref{app:lemma1}. The first lemma highlights a classical issue of under-specification in the Plackett-Luce model family \cite{plackett1975analysis}. Owing to this ambiguity, it is often necessary to introduce further identifiability conditions in order to guarantee meaningful maximum likelihood estimates as given in Eq. \ref{eq:reward_model} \cite{bong2022generalized}. The second lemma clarifies that all reward functions within the same class yield identical optimal policies, implying that our primary aim is to identify any representative from the correct equivalence class. The following theorem, proven in Appendix~\ref{app:thm1}, formalizes this intuition:

\begin{equation}\label{eq:main_eq}
    r(x,y) =\beta \log \frac{\pi_r(y\mid x)}{\pi_\text{ref}(y\mid x)} + \beta \log Z(x).
\end{equation}

\begin{equation}\label{eq:op_policy}
    \pi_r(y\mid x) = \frac{1}{Z(x)}\pi_\text{ref}(y\mid x)\exp\left(\frac{1}{\beta}r(x, y)\right),
\end{equation}

For the single-turn dialogue setting, our evaluation encompasses various approaches on a subset of the Anthropic HH dataset \citep{bai2022training} test split, specifically focusing on exchanges involving just one round between human and assistant. To assess GPT-4, we use the human-preferred completions in the test set as the ground truth and calculate win rates across methods. Due to the absence of a canonical SFT baseline for this benchmark, we begin with the pre-trained Pythia-2.8B model. We employ Preferred-FT to fine-tune a reference model using the selected completions, ensuring that generated completions remain close to the target distribution, followed by additional training with DPO. Our analysis further includes a comparison with the strongest output from 128 Preferred-FT generations (having observed in our experiments that the Best of $N$ metric saturates at $N=128$; see Appendix Figure~\ref{fig:best-of-n}), and a two-shot prompt version of the unmodified Pythia-2.8B model. We observe that DPO matches or exceeds the highest win rates across optimal sampling temperatures for each baseline. We also include an RLHF system, trained with PPO on the Anthropic HH dataset \footnote{\url{https://huggingface.co/reciprocate/ppo_hh_pythia-6B}}, as implemented by a reputable group \footnote{\url{https://github.com/CarperAI/trlx/tree/main/examples/hh}}. Despite extensive prompt and temperature tuning, this RLHF model does not surpass the base Pythia-2.8B's performance. Drawing on our findings from TL;DR as well as the shared reward maximization in both methods, we treat the Best of 128 baseline as an approximate indicator of PPO-level performance. In summary, DPO stands out as the only method that, while remaining computationally efficient, consistently outperforms the preferred completion benchmark on the Anthropic HH dataset, offering comparable or superior results relative to the resource-intensive Best of 128 approach. As depicted in Figure~\ref{fig:dialogue-main}, DPO also achieves its optimal performance after relatively few updates.

Given a fixed dataset of pairwise preferences $\mathcal{D}=\bigl\{x^{(i)}, y_w^{(i)}, y_l^{(i)}\bigr\}_{i=1}^N$ sampled from the underlying distribution $p^*$, we can represent the reward function as $r_{\phi}(x, y)$ and learn the associated parameters through maximum likelihood estimation. When formulating this as a binary classification problem, the associated negative log-likelihood loss is:

To further investigate how PPO and DPO handle distributional changes, we assessed their respective policies—trained on Reddit TL;DR summarization—on an out-of-distribution task using the test set of the CNN/DailyMail dataset \citep{nallapati-etal-2016-abstractive}, which consists of news articles. For this evaluation, we applied the optimal sampling temperatures identified in the TL;DR domain (0 and 0.25). Results, summarized in Table~\ref{tab:ood}, show that DPO continues to surpass the PPO policy by a wide margin on this new domain. Evaluation was conducted using the GPT-4 win rate metric, with GPT-4 prompted identically to the Reddit TL;DR assessment except substituting “news article” for “forum post.” These findings suggest that DPO-trained policies not only perform robustly in distributional shifts but also exhibit generalization capabilities on par with PPO policies, despite not leveraging the extra unlabeled Reddit TL;DR prompts utilized by PPO.

\begin{sproof}
    Take any reward function $r(x, y)$, which induces its corresponding optimal policy $\pi_r(y \mid x)$ as described by Eq. \ref{eq:op_policy}. We will demonstrate that a function from the equivalence class of $r$ can be constructed using the aforementioned reparameterization. Define the projection operator $f$ as  
\begin{equation}
    f(r; \pi_\text{ref}, \beta)(x, y) = r(x, y) - \beta\log\sum_{y}\pi_\text{ref}(y\mid x)\exp\left(\frac{1}{\beta}r(x, y)\right)
\end{equation}
Here, $f$ adjusts the reward by subtracting the log partition function (computed with $\pi_\text{ref}$) from $r(x, y)$. Because this normalization term is a function solely of $x$, $f(r; \pi_\text{ref}, \beta)(x, y)$ remains in the equivalence class of $r(x, y)$. Substituting $r$ using the right-hand side of Eq.~\ref{eq:main_eq} (which holds generally for any reward), yields $f(r; \pi_\text{ref}, \beta)(x, y) = \beta \log \frac{\pi_r(y\mid x)}{\pi_\text{ref}(y\mid x)}$. Thus, this projection provides a specific member of the equivalence class of $r$ in the desired form, so our proposed reparameterization is without loss of generality.
\end{sproof}

Having established a preference likelihood expressed in terms of the optimal policy instead of the reward, we can now define a maximum likelihood estimation objective for a parameterized policy $\pi_\theta$. In a manner akin to the reward modeling framework (see Eq.~\ref{eq:reward_model}), our objective for optimizing the policy is given by:

\begin{equation}\label{eq:AC}
    \max_{\pi_{\theta}}\mathbb{E}_{\pi_{\theta}(y\mid x)}\bigg[\underbrace{r_{\phi}(x, y) -\beta\log\sum_{y}\pi_\text{ref}(y\mid x)\exp\left(\frac{1}{\beta}r_{\phi}(x, y)\right)}_{f(r_{\phi}, \pi_\text{ref}, \beta)} - \underbrace{\beta\log\frac{\pi_{\theta}(y\mid x)}{\pi_\text{ref}(y\mid x)}}_{\text{KL}}\bigg]
\end{equation}

\subsection{Generalization to a new input distribution}

\begin{lemma}\label{lemma:same_policy}
    Any two reward functions that are members of the same equivalence class result in the same optimal policy when considered in the constrained RL setting.
\end{lemma}


\begin{abstract}
Large-scale unsupervised language models (LMs) acquire extensive world knowledge and some reasoning abilities; however, exerting fine-grained control over their outputs remains a challenge, primarily due to the inherently unsupervised nature of their training processes. To improve model steerability, common approaches involve gathering human judgments about the comparative quality of generated outputs, followed by fine-tuning the unsupervised LM to better align with these preferences—often through reinforcement learning from human feedback (RLHF). Yet, RLHF workflows are typically intricate and sometimes unstable, requiring an initial stage where a reward model is trained to capture human preferences, and a subsequent reinforcement learning phase where the original LM is updated to maximize the reward model’s output while maintaining proximity to its initial parameters. In this work, we propose an alternative reward model parameterization within RLHF that allows the derivation of the optimal policy in closed form. This reformulation enables us to address the conventional RLHF challenge by minimizing a straightforward classification loss. We introduce an algorithm, termed \textit{Direct Preference Optimization} (DPO), that is robust, efficient, and computationally streamlined—removing the necessity for language model sampling during fine-tuning and extensive hyperparameter optimization. Empirical results demonstrate that DPO fine-tunes LMs to human preferences as well as or better than prevailing methods. Specifically, DPO surpasses PPO-based RLHF in sentiment control and achieves comparable or superior quality in summarization and single-turn dialogue tasks, all while being much easier to implement and train.
\end{abstract}

\section{Experiments}
This section presents an empirical assessment of DPO’s capability to train policies using preference data. We begin by investigating, within a controlled text generation environment, how DPO balances the goals of reward maximization and KL-divergence minimization with respect to the reference policy in comparison to standard preference optimization techniques such as PPO. Following this, we explore DPO’s effectiveness on more complex settings, including larger model architectures and more challenging RLHF benchmarks, like summarization and dialogue. Our findings indicate that, with minimal hyperparameter adjustment, DPO consistently matches or surpasses the performance of robust baselines—such as PPO-based RLHF—as well as methods that select the best among $N$ sampled trajectories according to a learned reward model. The details of our experimental protocol are provided before the results, and further specifics are available in Appendix~\ref{app:exp_details}.

where $\sigma$ denotes the logistic sigmoid function. For language models, it is typical to initialize $r_{\phi}(x, y)$ with weights from the SFT model $\pi^\text{SFT}(y \mid x)$, extending it by attaching a linear head atop the final transformer layer to yield a scalar-valued reward prediction \cite{ziegler2020finetuning}. To mitigate high variance in the reward output, existing literature often normalizes the rewards such that $\mathbb{E}_{x,y\sim \mathcal{D}}\left[r_\phi(x, y)\right] = 0$ over the data distribution.

\begin{equation}\label{eq:reward_model}
    \mathcal{L}_R(r_{\phi}, \mathcal{D}) = -\mathbb{E}_{(x, y_w, y_l)\sim \mathcal{D}}\bigl[\log \sigma(r_{\phi}(x, y_w)- r_{\phi}(x, y_l))\bigr]
\end{equation}

\section{Theoretical Analysis of DPO} In this section, we present an in-depth interpretation of the DPO approach, establish its theoretical foundation, and connect DPO’s strengths to limitations inherent in actor-critic methods for RLHF, such as PPO~\cite{schulman2017proximal}.

\label{sec:theory}

\section{Discussion}
Preference-based learning presents a robust and extensible approach for developing language models that are both competent and aligned with human values. In this work, we proposed DPO, a straightforward training framework that allows for the training of language models from preference data without relying on reinforcement learning. Instead of adapting the preference learning problem to fit conventional reinforcement learning methodologies, DPO establishes a direct relationship between language model policies and reward structures. This connection enables training the model to conform to human preferences using a basic cross-entropy loss, circumventing the need for reinforcement learning and without any loss in generality. Remarkably, DPO matches or outperforms prevalent RLHF approaches, such as those leveraging PPO, with minimal hyperparameter adjustment; this substantially lowers the entry barrier for training language models based on human preference signals.

Acknowledging the computational and methodological obstacles of applying reinforcement learning at scale for language model fine-tuning, we aim to devise a streamlined alternative that directly optimizes the policy with respect to preference data. Departing from the standard RLHF pipeline, which involves reward learning followed by RL optimization, our method capitalizes on a special form of reward model that allows its optimal policy to be derived analytically, obviating the need for a separate RL loop. As detailed below, our primary innovation is to exploit a closed-form correspondence between reward functions and their corresponding optimal policies, thereby recasting a loss over rewards into an equivalent loss over policies. This change-of-variables framework enables us to bypass explicit reward model fitting, yet remains fully compatible with

\section{Introduction}
Large-scale unsupervised language models (LMs) trained on massive corpora have demonstrated unexpected capabilities~\citep{chowdhery2022palm, brown2020language, touvron2023llama,bubeck2023sparks}. Nonetheless, the training data for these models comprises human-generated content that reflects a broad spectrum of intentions, priorities, and expertise. Not all of these qualities are desirable for imitation; for instance, while an AI coding assistant should \textit{recognize} common coding errors to provide corrections, we prefer that its code generation is influenced primarily by the (sometimes infrequent) high-caliber coding examples in the dataset. In a similar vein, it is important for a language model to be \textit{informed} about widely held misconceptions—such as those believed by half the population—yet we would not want the model to affirm these misconceptions in half of its relevant responses. Thus, the task of constraining the model’s \emph{actual outputs and conduct} based on its vast \textit{underlying knowledge and skills} is fundamental for constructing AI systems that are reliable, effective, and manageable \citep{ouyang2022training}. Although prevailing techniques direct LMs to align with human values using reinforcement learning (RL), this work demonstrates that the RL-derived objective can be optimized precisely using a straightforward binary cross-entropy loss, leading to a significant simplification of the preference learning framework.

\textbf{Evaluation.} We employ two complementary strategies for evaluating our models. To assess each algorithm's capability in optimizing the constrained reward maximization objective, we use the controlled sentiment generation context, where we plot each method's tradeoff between achieved reward and KL-divergence from a reference policy. This evaluation is possible because we possess access to the true reward function (implemented as a sentiment classifier). In contrast, in practical applications where the true reward function is unknown, we adopt a different approach: we compute the algorithms’ \textit{win rate} relative to a baseline policy. Here, GPT-4 acts as a stand-in for human evaluators, judging summary quality and helpfulness for summarization and single-turn dialogue tasks, respectively. For summarization, our baseline comprises reference summaries from the test set; for dialogue, it consists of the preferred responses from test data. While recent work indicates that language models may offer more consistent automated evaluation than traditional metrics \citep{Chen2023ExploringTU}, we also conduct a human evaluation to support our reliance on GPT-4 as an evaluator, as detailed in Sec.~\ref{sec:human-judgments}. Our analysis reveals that GPT-4's assessments are highly consistent with human judgments, with human-GPT-4 agreement rates generally on par with, or exceeding, agreement between different human annotators.

Large self-supervised language models have demonstrated the capacity to solve certain tasks without prior examples (zero-shot) \citep{radford2019language} or with minimal prompting (few-shot) \citep{gpt3,megatron,chowdhery2022palm}. Nonetheless, their effectiveness on downstream applications and their alignment with user intentions are greatly enhanced by fine-tuning on curated datasets containing instructions and human-generated responses \citep{mishra-etal-2022-cross,sanh2022multitask,chung2022scaling,thoppilan2022lamda}. This process, known as 'instruction-tuning,' allows large language models (LLMs) to extrapolate beyond the scope of their training instructions and generally boosts their practical utility \citep{chung2022scaling}. In practice, while instruction-tuning has been quite successful, obtaining \textit{relative} human evaluations of output quality is typically simpler than acquiring expert-annotated examples. As a result, subsequent research has adapted LLMs using datasets reflecting human preferences, leading to improvements in tasks like translation \citep{kreutzer-etal-2018-reliability}, summarization \citep{stiennon2022learning,ziegler2020finetuning}, story generation \citep{ziegler2020finetuning}, and following instructions \citep{ouyang2022training,ramamurthy2023is}. These approaches involve first training a neural network-based reward model that reflects human preference data, often under frameworks like the Bradley-Terry model \citep{bradley1952rankanalysis}, followed by refining the language model to maximize this reward using reinforcement learning strategies such as REINFORCE \citep{williams1992reinforce}, proximal policy optimization (PPO; \cite{schulman2017proximal}), or their modifications \citep{ramamurthy2023is}. Parallel research also employs instruction-tuned LLMs, further enhanced with human feedback, to autonomously generate synthetic preference datasets targeting specific attributes like safety or harmlessness \citep{bai2022constitutional}. Here, minimal human oversight is provided in the form of text-based rubrics for LLM evaluation. Collectively, these methodologies unify two areas of research: reinforcement learning-based language model training for various objectives~\citep{Ranzato2015SequenceLT,paulus2018a,wu2018learning}, and broader strategies for preference-based learning from humans \citep{christiano2017deep,kupcsik2018learning}. Although optimizing LLMs with relative human feedback holds considerable promise, reinforcement learning-based fine-tuning for large models continues to be a substantial practical hurdle. The present work introduces a theoretically grounded method for optimizing models based on relative preferences that does not require reinforcement learning.

It can be readily verified that this defines an equivalence relation, thereby dividing the space of reward functions into distinct equivalence classes. The following two lemmas then hold:

\textbf{Derivation of the DPO Objective.} We begin with the standard RL objective as given in Eq.~\ref{eq:RL}, utilizing a general reward function $r$. As established in previous work~\citep{peters2007reinforcement, peng2019advantage, korbak2022reinforcement, go2023aligning}, the solution that optimally satisfies the KL-constrained reward maximization problem (see Eq.~\ref{eq:RL}) can be explicitly written as:

In general, current approaches impart desirable characteristics to language models by leveraging curated datasets of human preferences that reflect behaviors deemed safe and beneficial. This stage of preference alignment follows an initial unsupervised pre-training phase on extensive textual data. Although supervised fine-tuning on exemplary human-provided answers is the most direct method for preference alignment, the predominant and most effective paradigm is reinforcement learning from human or AI feedback (RLHF/RLAIF; \citep{christiano2017deep,bai2022constitutional}). In RLHF, a reward model is learned from annotated human preference data, and reinforcement learning is then used to refine the LM policy to favor outputs that receive high reward, while discouraging excessive divergence from the pre-trained model. While RLHF has led to models with remarkable conversational and programming prowess, its methodology is considerably more involved than supervised fine-tuning: it requires training several LMs and continuously sampling from the model’s policy throughout training, thereby imposing substantial computational demands.

\begin{equation}\label{eq:objective}
    p^*(y_1\succ y_2 \mid x)=\frac{1}{1 + \exp\left(\beta \log \frac{\pi^*(y_2\mid x)}{\pi_\text{ref}(y_2\mid x)} - \beta \log \frac{\pi^*(y_1\mid x)}{\pi_\text{ref}(y_1\mid x)}\right)}
\end{equation}

Here, $\hat{r}_\theta(x, y) = \beta \log \frac{\pi_\theta(y \mid x)}{\pi_\text{ref}(y \mid x)}$ specifies the implicit reward assigned by the current language model $\pi_\theta$ relative to a reference policy $\pi_\text{ref}$ (see additional discussion in Section~\ref{sec:theory}). Conceptually, the gradient update for $\mathcal{L}_\text{DPO}$ strengthens the probability of preferred responses $y_w$, while reducing the probability of dispreferred alternatives $y_l$. Critically, the contribution of each training example is modulated by how much the implicit reward model $\hat{r}_\theta$ favors the less desirable completion, modulated by $\beta$—in other words, the degree to which the reward model's ranking deviates from the provided preference, factoring in the strength of the KL regularization. Empirical results highlight the necessity of this weighting scheme, as a simpler variant that omits the weighting factor can lead to pathological behaviors in the language model (see Appendix Table~\ref{tab:unlikelihood_generations}).

\section{Preliminaries}\label{section:prelims}

To assess the trustworthiness of GPT-4’s evaluations, we ran a human annotation study based on TL;DR summarization results and tested two GPT-4 prompts. The \textbf{GPT-4 (S)} (simple) prompt requests a judgment on which summary better captures essential information, while the \textbf{GPT-4 (C)} (concise) prompt adds an explicit instruction to weigh conciseness as well, motivated by observations that GPT-4 (under the \textbf{GPT-4 (S)} prompt) tends to select summaries that are longer and more redundant than those favored by human evaluators. Full prompt texts are included in Appendix~\ref{app:prompts}. We carried out three types of comparisons—one using the highest-performing model (DPO, temp. 0.25), one with the lowest (PPO, temp. 1.0), and one intermediary (SFT, temp. 0.25)—to span a broad spectrum of summary qualities. Each was evaluated against the best-performing PPO policy sampled greedily. Our results reveal that, regardless of prompt, GPT-4’s judgments align with human preferences at a rate comparable to inter-human agreement, supporting the use of GPT-4 as a stand-in for human raters (for DPO and PPO-1, we collected multiple human annotations due to limited available raters). The \textbf{GPT-4 (C)} prompt, in particular, yields win rates most consistent with human perspectives and is thus used in Section~\ref{sec:dpo-real-datasets} for principal results. Further methodological details, including the user interface and identities of human participants, are provided in Appendix~\ref{app:human-study}.

\begin{equation}\label{eq:lag_p}
     \sum_{y}\underbrace{\pi_\text{ref}(y\mid x)\exp\left(\frac{1}{\beta}r(x, y)\right)}_{=\pi(y\mid x)\text{, using Thm.~\ref{thm:main} reparam.}} = 1,
\end{equation}

\section{Direct Preference Optimization}\label{sec:DPO}

A detailed derivation is available in Appendix~\ref{app:derivation2}. While Eq.~\ref{eq:objective} leverages the Bradley-Terry approach, this derivation naturally extends to the more general Plackett-Luce models~\citep{plackett1975analysis, luce2012individual}, as outlined in Appendix~\ref{app:plackett_luce_models}.

where the partition function is defined by $Z(x) =\sum_{y}\pi_\text{ref}(y\mid x)\exp\left(\frac{1}{\beta}r(x, y)\right)$. For a full step-by-step derivation, refer to Appendix \ref{app:derivation1}. Even if we substitute the maximum likelihood estimator $r_{\phi}$ for the true reward function $r^*$, computing the partition function $Z(x)$ remains computationally demanding~\citep{korbak2022reinforcement, go2023aligning}, which limits the practical usability of this representation. However, by manipulating Eq.~\ref{eq:op_policy}, we can express the reward function in terms of the optimal policy $\pi_r$, the reference policy $\pi_\text{ref}$, and the partition function $Z(\cdot)$. Specifically, by applying the logarithm to both sides of Eq.~\ref{eq:op_policy} and simplifying, we arrive at:

\begin{lemma}\label{lemma:same_prefrence} Within the Plackett-Luce framework, including the Bradley-Terry model as a special case, reward functions belonging to the same equivalence class yield identical preference distributions.
\end{lemma}

which ensures that $\pi(y\mid x)$ constitutes a legitimate probability distribution (all probabilities are non-negative and sum to one). Moreover, referencing Eq.~\ref{eq:op_policy}, we note that Eq.~\ref{eq:lag_p} corresponds to the partition function for the optimal policy derived from the reward function $r(x, y)$. The fundamental insight behind the DPO algorithm is the realization that we can enforce suitable constraints on the otherwise underdetermined Plackett-Luce family of preference models (with the Bradley-Terry model as a prominent example) to maintain the expressive power of representable reward models. At the same time, these constraints ensure that the optimal policy described by Eq. \ref{eq:op_policy} remains analytically manageable for any prompt $x$.

\begin{definition}
Two reward functions $r(x, y)$ and $r'(x, y)$ are defined as equivalent if and only if there exists a function $f$ such that $r(x, y)-r'(x, y) = f(x)$.     
\end{definition}

\textbf{Methods.} Beyond DPO, our study benchmarks several prominent approaches for training language models to match human preferences. For baseline prompting, we consider zero-shot prompts to \textbf{GPT-J} \citep{gpt-j} in the summarization domain and 2-shot prompts to \textbf{Pythia-2.8B} \citep{biderman2023pythia} in dialogue. We also include the \textbf{SFT} model, as well as \textbf{Preferred-FT}—a supervised model fine-tuned on selected completions $y_w$ generated either by the SFT model (in both controlled sentiment and summarization settings) or a general-purpose language model (in dialogue). Another comparable approach is the \textbf{Unlikelihood} method~\citep{welleck2019neural}, which optimizes the model to increase the likelihood of $y_w$ while explicitly decreasing the likelihood of $y_l$, with an optional scaling parameter $\alpha\in[0,1]$ on the unlikelihood loss. We further examine \textbf{PPO} \citep{schulman2017proximal}, trained on rewards inferred from preference data, and \textbf{PPO-GT}, an oracle variant that leverages the known ground-truth reward in the controlled sentiment experiments. Specifically, for sentiment tasks, we apply two PPO-GT implementations: a standard version \cite{leandro_von_werra_2023_7790115}, and a modified one that normalizes the reward and fine-tunes hyperparameters for enhanced results (these adjustments are likewise applied to the standard PPO with learned rewards). Lastly, we evaluate the \textbf{Best of $N$} baseline, wherein $N$ samples are drawn from the SFT (or Preferred-FT for dialogue), and the sample with the highest score under a reward model trained on preference data is chosen. Although this strategy tends to yield high performance by separating reward modeling from PPO optimization, its requirement to generate $N$ completions for every input makes it computationally burdensome, even for modest $N$ values, at inference time.